\documentclass[11pt]{scrartcl}

\usepackage{amsmath,amsthm,amssymb,mathtools}
\usepackage{hyperref}
\usepackage{lastpage}
\usepackage{xcolor}
\usepackage{multicol}

\theoremstyle{plain} 
\newtheorem{theorem}{Theorem}[section]
\newtheorem{lemma}[theorem]{Lemma}
\newtheorem*{question}{Question}
\theoremstyle{definition}
\newtheorem*{definition}{Definition} 
\theoremstyle{remark}
\newtheorem{remark}{Remark}

\usepackage[none]{hyphenat}
\setlength{\parindent}{0pt}
\setlength{\emergencystretch}{3em}

\renewcommand\labelitemi{$\vcenter{\hbox{\tiny$\bullet$}}$} 

\usepackage[margin=1.5in]{geometry}
\usepackage{fancyhdr}
\pagestyle{fancy}
\lhead{ESC194: Calculus I} 
\rhead{\thepage}
\cfoot{}

\title{ESC194: Calculus I \\ Exam Review}
\author{Brendan Tang}
\date{Fall 2026}

\begin{document}

\maketitle
\thispagestyle{empty}

\begin{abstract}
    \centering
    3 months of Calculus I in \pageref*{LastPage} pages.
\end{abstract}

\tableofcontents
\newpage

\section{Limits and Continuity}
\subsection{The Precise Definition of a Limit}

The precise definition of a limit is as follows.

\begin{definition}
    If $\forall\epsilon>0$, $\exists\delta>0$ such that $\lvert x-a \rvert<\delta\implies
    \lvert f(x)-L \rvert<\epsilon$, then \[ \lim_{x \to a}f(x)=L. \]
\end{definition}

For one-sided limits, $x$ is only approaching $a$ from one side. Thus, we define the left-handed
limit as follows.

\begin{definition}
    If $\forall\epsilon>0$, $\exists\delta>0$ such that $a-\delta<x<a\implies \lvert f(x)-L \rvert
    <\epsilon$, then \[ \lim_{ x \to a^- }f(x)=L. \]
\end{definition}

As a number becomes arbitrarily large, it eventually surpasses every real number. Thus, we define
limits to or approaching infinity as follows.

\begin{definition}
    If $\forall M>0$, $\exists\delta>0$ such that $0<\lvert x-a \rvert<\delta\implies f(x)>M$, then
    \[ \lim_{ x \to a }f(x) \to \infty. \]
    Conversely, if $\forall\epsilon>0$, $\exists M>0$ such that $x>M\implies \lvert f(x)-L \rvert
    <\epsilon$, then \[ \lim_{ x \to \infty }f(x)=L. \]
\end{definition}

\begin{question}
    How can we rigorously define right-handed and negative infinite limits? What about limits 
    involving both concepts?
\end{question}

Here are the general steps to a rigorous proof.

\begin{enumerate}
    \item Given $\epsilon>0$, find $\delta>0$. Start with $f(x)-L<\epsilon$ and manipulate
    it into the form $\lvert x-a \rvert<\delta_{0}$. This step may require bounding $\delta\leq n$.
    \item \begin{proof} 
        Given $\epsilon$, let $\delta=min(n,\delta_{0})$. Then, show that $\lvert x-a 
        \rvert<\delta\implies \lvert f(x)-L \rvert<\epsilon$, as desired. 
    \end{proof}
\end{enumerate}

\subsection{Limit Theorems}

Here are some important limit theorems.

\begin{theorem}
    If $\lim_{ x \to a }f(x)$ and $\lim_{ x \to a }g(x)$ exist, then
    \begin{itemize}
        \item (Constant LT) $\lim_{ x \to a }c=c$
        \item (Sum LT) $\lim_{ x \to a }[f(x)+g(x)]=\lim_{ x \to a }f(x)+\lim_{ x \to a }g(x)$
        \item (Product LT) $\lim_{ x \to a }[f(x)\cdot g(x)]=\lim_{ x \to a }f(x)\cdot 
        \lim_{ x \to a }g(x)$
        \item (Constant multiple LT) $\lim_{ x \to a }cf(x)=c\lim_{ x \to a }f(x)$
        \item (Power LT) $\lim_{ x \to a }[f(x)]^n=[\lim_{ x \to a }f(x)]^n$
        \item (Root LT) $\lim_{ x \to a }\sqrt[n]{ f(x) }=\sqrt[n]{ \lim_{ x \to a }f(x) }$
    \end{itemize}
\end{theorem}

\begin{theorem}[Direct substitution]
    If $f$ is continuous at $a$, then \[ \lim_{ x \to a } f(x)=a. \]
\end{theorem}

The following theorem is often used to evaluate trigonometric limits.

\begin{theorem}[Squeeze theorem]
    If $g(x)<f(x)<h(x)$ and $\lim_{ x \to a }g(x)=\lim_{ x \to a }h(x)=L$, then 
    \[ \lim_{ x \to a } f(x)=L. \]
\end{theorem}

When solving difficult trigonometric limits, it may help to find $\frac{{\sin x}}{x}$ terms.

\subsection{Continuity Theorems}

Here is the definition for continuity.

\begin{definition}
    $f$ is continuous at a point $a$ if $\lim_{ x \to a }f(x)=f(a)$. 
\end{definition}

\begin{question} 
    How can we define continuity from the side? What about continuity on an open or closed interval?
\end{question}

Here are some continuity theorems.

\begin{theorem}
    If $f$ and $g$ are continuous at $a$, then \[ f+g,\ fg,\ cf \] are also continuous at $a$.
\end{theorem}

\begin{theorem}
    The following types of functions are continuous on their domains 
    \begin{multicols}{3}
        \begin{itemize}
            \item polynomial
            \item exponential
            \item rational
            \item logarithmic
            \item root
            \item trigonometric
        \end{itemize}
    \end{multicols}
\end{theorem}

In fact, all elementary functions are continuous on their domains. 
\newpage

\section{Derivatives}
\subsection{Derivative Theorems}

We begin with the definition of the derivative.

\begin{definition}
    The derivative of $f$ is defined as \[ f'(x)\coloneqq \lim_{ h \to 0 } \frac{{f(x+h)-f(x)}}{h}. \]
\end{definition}

Now we define differentiability. 

\begin{definition}
    $f$ is differentiable at a point $a$ if $f'(a)$ exists.
\end{definition}

\begin{question} 
    How can we define differentiability from the side? What about differentiability on an open or
    closed interval?
\end{question}

From the definition of the derivative, we see that differentiability at $a$ requires $\lim_{ h \to 0 } 
[f(a+h)-f(a)]=0$ (why?). This implies that $f$ is continuous at $a$. Thus, we arrive at the following theorem.

\begin{theorem}
    If $f$ is differentiable at $a$, then $f$ is continuous at $a$.
\end{theorem}

Be careful, as this theorem does not hold for the converse. That is, continuity at $a$ does not imply
differentiability at $a$.

\end{document}